\documentclass[10pt]{article}
\usepackage[letterpaper,margin=0.72in]{geometry}
\usepackage[T1]{fontenc}
\usepackage{lmodern,microtype,xcolor,enumitem,hyperref}
\definecolor{navy}{HTML}{08111F}
\definecolor{blue}{HTML}{2677FF}
\definecolor{muted}{HTML}{5D6B80}
\hypersetup{colorlinks=true,urlcolor=blue}
\setlist[itemize]{leftmargin=1.2em,itemsep=2pt,topsep=3pt}
\pagestyle{plain}
\begin{document}
{\color{blue}\small\bfseries SUMMER ENGINEERING PROJECT --- 2026}
\vspace{0.4em}
{\color{navy}\LARGE\bfseries Polargraph Chalk Plotter}\\
{\color{muted}\large Cable-driven mechatronics for large-format wall drawing}

\section*{Project overview}
A cable-driven drawing machine that converts planar artwork into coordinated motor motion and deposits chalk on a vertical surface. The project combines inverse kinematics, step generation, end-effector design, calibration, and structural iteration.

\section*{Engineering problem}
Large-format plotting is mechanically simple in concept but sensitive to cable stretch, spool geometry, wall friction, missed steps, and chalk contact force.

\section*{System architecture}
\begin{itemize}
\item Two upper motor-spool assemblies control cable lengths to position the carriage.
\item A chalk carriage provides compliance, repeatable contact, and a controlled lift mechanism.
\item A motion planner converts vector paths into short line segments and synchronized steps.
\item A calibration routine maps commanded cable length to measured wall coordinates.
\end{itemize}

\section*{Engineering focus}
\begin{itemize}
\item Inverse kinematics from target (x, y) position to left and right cable lengths.
\item Step-density and acceleration limits selected to reduce oscillation and missed steps.
\item Spool-radius compensation and homing references for repeatability.
\end{itemize}

\section*{Verification plan}
\begin{itemize}
\item Plot a calibration grid and measure scale, skew, closure, and repeatability errors.
\item Test circles, diagonals, corners, and long strokes at increasing speed.
\item Log failure modes such as cable slip, chalk chatter, and carriage tilt.
\end{itemize}

\section*{Status and evidence boundary}
Summer 2026 mechatronics prototype brief; quantitative accuracy belongs in a later test report.

\vfill
\small\color{muted} V. Narayan Kushwaha --- Summer Engineering Portfolio 2026
\end{document}
